Floating platforms let wind turbines be placed in deep water, but the platform moves, and its
motion couples with the rotor in ways that a controller designed for a fixed tower does not
anticipate. This dissertation treats the platform geometry and the feedback controller as one
design problem. It formulates a co-design optimization whose variables include ballast, column
spacing and controller gains, subject to fatigue and pitch-motion limits. Solving it for a
15-megawatt reference turbine yields a platform 12\% lighter than the sequentially designed
baseline at equal fatigue damage.
